#CSS Learning notes
In CSS the order is important, the last rule will override the previous ones if they have the same specificity, for this reason is named "Cascading Style Sheets".

## Color
Colors can be defined in several ways, the most common are:
- Named colors: `red`, `blue`, `green`, etc.
- Hexadecimal: `#RRGGBB` where RR, GG, and BB are two-digit hexadecimal numbers representing the red, green, and blue components of the color. And, we can use two additional digits for alpha (transparency) in the format `#RRGGBBAA`.
- RGB: `rgb(R, G, B)` where R, G, and B are integers between 0 and 255 representing the red, green, and blue components of the color. We can also use `rgba(R, G, B, A)` where A is a number between 0 and 1 representing the alpha (transparency) component.
  Or use 'rgb(R G B / A)' where R, G, and B are integers between 0 and 255 representing the red, green, and blue components of the color, and A is a number between 0 and 1 representing the alpha (transparency) component.
- HSL: `hsl(H, S\%, L\%)` where H is an integer between 0 and 360 representing the hue, S is a percentage representing the saturation, and L is a percentage representing the lightness. We can also use `hsla(H, S\%, L\%, A)` where A is a number between 0 and 1 representing the alpha (transparency) component.
- oklch: `oklch(L C H)` where L is a number between 0 and 1 representing the lightness, C is a number representing the chroma, and H is an integer between 0 and 360 representing the hue. We can also use `oklch(L C H / A)` where A is a number between 0 and 1 representing the alpha (transparency) component. 

## Special values
### iherit
The `inherit` value allows an element to inherit a property value from its parent element.
### initial
The `initial` value resets a property to its default value as defined in the CSS specification.
### unset
The `unset` value resets a property to its inherited value if it inherits from its parent, or to its initial value if it does not inherit.
### revert
The `revert` value resets a property to the value established by the user-agent stylesheet


## Pseudo-classes
### hover
The `:hover` pseudo-class applies when the user designates an element with a pointing device, but does not activate it. 
For example, when the mouse pointer is over an element, the `:hover` pseudo-class can be used to change some attributes of the element, such as its background color or text color.

### active
The `:active` pseudo-class applies when an element is being activated by the user.
For example, when a user clicks on a button, the `:active` pseudo-class can be used to change the appearance of the button while it is being clicked.

### focus
The `:focus` pseudo-class applies when an element has received focus, either from the user selecting it with the mouse or by navigating to it using the keyboard.
For example, when a user clicks on an input field or navigates to it using the keyboard, the `:focus` pseudo-class can be used to change the appearance of the input field to indicate that it is active and ready for user input.

### focus-visible
The `:focus-visible` pseudo-class applies when an element has received focus and the user agent determines that it should be visibly focused.

### focus-within
The `:focus-within` pseudo-class applies to an element if it or any of its descendants are focused.
This is useful for styling parent elements when any of their child elements receive focus, such as highlighting a form group when an input field within it is focused.

### target
The `:target` pseudo-class applies to an element that is the target of a fragment identifier in the URL.
For example, if the URL is `http://example.com/page#section1`, the element with the id `section1` will match the `:target` pseudo-class, allowing you to style it differently when it is the target of the URL.

### first-child
The `:first-child` pseudo-class applies to an element that is the first child of its parent.

### last-child
The `:last-child` pseudo-class applies to an element that is the last child of its parent.


## Important Things
### border vs. outline
The `border` property adds a border around an element, while the `outline` property adds a line around an element that does not take up space and does not affect the layout of the page. 
The `outline` is often used for accessibility purposes, such as indicating focus on interactive elements.